% !TEX program = xelatex
% ============================================================================
% 《人工智能基础》大作业技术报告 —— 仓库穿梭无人车（AGV）智能调度
% 编译方式：tectonic main.tex  或  Overleaf 选择 XeLaTeX 编译
% ============================================================================
\documentclass[UTF8, a4paper, zihao=-4, fontset=none]{ctexart}

% ============================================================================
% 公共样式与跨平台排版配置（参考 NeurIPS 论文排版比例，纯黑正文字体 + 彩色配图）
% ============================================================================

% --- 页面尺寸与页边距（NeurIPS 单栏学术论文版芯比例）---
\usepackage[a4paper, top=2.4cm, bottom=2.4cm, left=2.7cm, right=2.7cm, headheight=14pt]{geometry}

% --- 跨平台中英文字体配置（自动适配 Windows / macOS / Linux & Overleaf）---
\usepackage{fontspec}
\IfFontExistsTF{Times New Roman}{
  \setmainfont{Times New Roman}
}{
  \IfFontExistsTF{TeX Gyre Termes}{
    \setmainfont{TeX Gyre Termes}
  }{}
}
\IfFontExistsTF{Arial}{
  \setsansfont{Arial}
}{
  \IfFontExistsTF{TeX Gyre Heros}{
    \setsansfont{TeX Gyre Heros}
  }{}
}
\IfFontExistsTF{Menlo}{
  \setmonofont[Scale=0.86]{Menlo}
}{
  \IfFontExistsTF{DejaVu Sans Mono}{
    \setmonofont[Scale=0.85]{DejaVu Sans Mono}
  }{}
}

% 正文中文字体：宋体（优先 SimSun -> Songti SC -> STSong -> FandolSong）
\IfFontExistsTF{SimSun}{
  \setCJKmainfont[AutoFakeBold=2.2, ItalicFont=KaiTi]{SimSun}
}{
  \IfFontExistsTF{Songti SC}{
    \setCJKmainfont[BoldFont=Songti SC Bold, ItalicFont=Kaiti SC]{Songti SC}
  }{
    \IfFontExistsTF{STSong}{
      \setCJKmainfont[AutoFakeBold=2.2, ItalicFont=Kaiti SC]{STSong}
    }{
      \setCJKmainfont[AutoFakeBold=2.2]{FandolSong}
    }
  }
}

% 无衬线中文字体：黑体
\IfFontExistsTF{SimHei}{
  \setCJKsansfont[AutoFakeBold=1.8]{SimHei}
}{
  \IfFontExistsTF{Heiti SC}{
    \setCJKsansfont[BoldFont=Heiti SC Medium]{Heiti SC}
  }{
    \IfFontExistsTF{STHeiti}{
      \setCJKsansfont[AutoFakeBold=1.8]{STHeiti}
    }{
      \setCJKsansfont[AutoFakeBold=1.8]{FandolHei}
    }
  }
}

% 等宽中文字体与楷体定义
\IfFontExistsTF{FangSong}{
  \setCJKmonofont{FangSong}
}{
  \IfFontExistsTF{Kaiti SC}{
    \setCJKmonofont{Kaiti SC}
  }{
    \setCJKmonofont{FandolFang}
  }
}
\IfFontExistsTF{KaiTi}{
  \newCJKfontfamily\kaiti{KaiTi}
}{
  \IfFontExistsTF{Kaiti SC}{
    \newCJKfontfamily\kaiti{Kaiti SC}
  }{
    \newCJKfontfamily\kaiti{FandolKai}
  }
}

% --- 数学、表格、图形与算法宏包 ---
\usepackage{amsmath, amssymb, amsthm}
\usepackage{booktabs}
\usepackage{tabularx}
\usepackage{multirow}
\usepackage{array}
\usepackage{graphicx}
\usepackage{xcolor}
\usepackage{float}
\usepackage{caption}
\usepackage{subcaption}
\usepackage{enumitem}
\usepackage{setspace}
\usepackage{tikz}
\usetikzlibrary{arrows.meta, positioning, shapes.geometric, fit, backgrounds, calc}
\usepackage[most]{tcolorbox}

% --- 参考文献与超链接 ---
\usepackage[numbers, sort&compress]{natbib}
\usepackage[hidelinks, bookmarksnumbered=true]{hyperref}

% --- 配图专用色彩定义（仅用于 TikZ 图形内部，标题与正文全部为纯黑）---
\definecolor{PrimaryNavy}{HTML}{1B365D}
\definecolor{AccentBlue}{HTML}{2B6CB0}
\definecolor{SoftBlueBg}{HTML}{EDF2F7}
\definecolor{AccentTeal}{HTML}{2C7A7B}
\definecolor{SoftTealBg}{HTML}{E6FFFA}
\definecolor{WarmAmber}{HTML}{B7791F}
\definecolor{SoftAmberBg}{HTML}{FEFCBF}
\definecolor{BorderGray}{HTML}{CBD5E0}
\definecolor{MutedText}{HTML}{4A5568}

% --- 紧凑学术行距与段落间距（对齐 NeurIPS 10pt 论文视觉密度）---
\linespread{1.18}
\setlength{\parskip}{3.5pt}
\setlength{\emergencystretch}{2em}
\setlist[itemize]{itemsep=1.5pt, topsep=3pt, parsep=0pt}
\setlist[enumerate]{itemsep=1.5pt, topsep=3pt, parsep=0pt}

% --- 章节标题格式（参考 NeurIPS 分级字号：纯黑加粗，紧凑克制）---
% 正文基准为 10pt (五号 10.5pt)：
% 一级标题 \section: 12pt 加粗
% 二级标题 \subsection: 10.5pt 加粗
% 三级标题 \subsubsection: 10pt 加粗
\ctexset{
  section = {
    format = \fontsize{12pt}{15pt}\selectfont\bfseries,
    name = {,},
    number = \arabic{section},
    aftername = \hspace{0.55em},
    beforeskip = 12pt plus 3pt minus 2pt,
    afterskip = 5pt plus 1pt minus 1pt,
  },
  subsection = {
    format = \fontsize{10.5pt}{13.5pt}\selectfont\bfseries,
    aftername = \hspace{0.45em},
    beforeskip = 9pt plus 2pt minus 1pt,
    afterskip = 4pt plus 1pt minus 1pt,
  },
  subsubsection = {
    format = \fontsize{10pt}{12.5pt}\selectfont\bfseries,
    aftername = \hspace{0.4em},
    beforeskip = 6pt plus 1pt minus 1pt,
    afterskip = 3pt plus 1pt minus 1pt,
  }
}

% --- 图表标题格式（纯黑紧凑 9pt）---
\captionsetup{
  font={footnotesize},
  labelfont={bf},
  labelsep=colon,
  skip=5pt
}

% --- 页眉与页脚严格规则 ---
% 第 1 页：顶部居中显示《人工智能基础》大作业 + 单横线（对应 P2）
% 第 2 页及以后：完全取消页眉文字与页眉横线（删去 P4），仅保留底部居中页码
\usepackage{fancyhdr}

\fancypagestyle{firstpage}{
  \fancyhf{}
  \fancyhead[C]{\fontsize{9.5pt}{11pt}\selectfont 《人工智能基础》大作业}
  \renewcommand{\headrulewidth}{0.6pt}
  \renewcommand{\footrulewidth}{0pt}
  \fancyfoot[C]{\small\thepage}
}

\fancypagestyle{mainbody}{
  \fancyhf{}
  \renewcommand{\headrulewidth}{0pt}
  \renewcommand{\footrulewidth}{0pt}
  \fancyfoot[C]{\small\thepage}
}

\pagestyle{mainbody}

% --- 队友留白占位框（纯黑文字配浅灰虚线框，标明对应队友提交模块）---
\newtcolorbox{teammatebox}[1]{
  enhanced,
  breakable,
  colback=black!2,
  colframe=black!45,
  coltitle=black,
  fonttitle=\bfseries\footnotesize,
  title={队友负责模块留白（待队友补充）：#1},
  boxrule=0.5pt,
  arc=2pt,
  left=8pt,
  right=8pt,
  top=5pt,
  bottom=5pt,
  fontupper=\footnotesize,
  attach boxed title to top left={xshift=8pt, yshift=-2mm},
  boxed title style={colback=white, colframe=black!45, boxrule=0.5pt, arc=2pt}
}

% --- 论文标题与双人居中署名区（参考上传 NeurIPS 论文标题与作者字号比例，纯黑）---
% 标题：上下双横线（上粗 2.5pt、下细 0.8pt），字号 16pt 加粗居中
% 学号与姓名：与 NeurIPS 作者栏一致采用 10pt 加粗居中，两行紧凑排列
\newcommand{\MakePaperTitle}[5]{%
  \thispagestyle{firstpage}%
  \begin{center}
    \vspace*{2pt}
    \noindent\rule{\textwidth}{2.4pt}\par
    \vspace{6pt}
    {\fontsize{15.5pt}{21pt}\selectfont\bfseries #1 \par}
    \vspace{4pt}
    \noindent\rule{\textwidth}{0.7pt}\par
    \vspace{9pt}
    {\fontsize{10pt}{14.5pt}\selectfont\bfseries
      学号：#2，姓名：#3 \\[2pt]
      学号：#4，姓名：#5 \par
    }
    \vspace{8pt}
  \end{center}
}

% --- 代码标识符行内等宽排版 ---
\newcommand{\code}[1]{\texttt{#1}}


% ============================================================================
% 组员信息配置（整个队伍仅 2 人，请直接修改下方四项）
% ============================================================================
\newcommand{\ReportTitle}{多任务并发下的仓库穿梭无人车（AGV）智能调度方法研究与实现}
\newcommand{\StudentOneID}{待填写学号1}
\newcommand{\StudentOneName}{待填写姓名1}
\newcommand{\StudentTwoID}{待填写学号2}
\newcommand{\StudentTwoName}{待填写姓名2}

\begin{document}

% 首页标题与双人居中署名（首页启用 firstpage 页眉，后续页面自动切换为无页眉 mainbody）
\MakePaperTitle{\ReportTitle}{\StudentOneID}{\StudentOneName}{\StudentTwoID}{\StudentTwoName}

% 摘要与关键词
% ============================================================================
% 摘要与关键词
% ============================================================================
\begin{center}
  \begin{minipage}{0.92\textwidth}
    \small
    \noindent{\sffamily\bfseries 摘\quad 要：}
    针对自动化立体仓库中多取货任务并发场景下的多穿梭无人车（AGV）智能调度问题，本文围绕“动态任务指派”与“无碰撞时空路径规划”的强耦合难题，设计并实现了一套集建模仿真、启发式调度优化与独立物理验解于一体的调度系统。系统采用“任务指派层—时空路径规划层—独立物理验解层”三层协同架构，在严格遵守货架层序解锁、窄巷道顶点与换边避碰、装卸占位及流关闭返库等物理约束的前提下，围绕通用多工况调度引擎（\code{src/agv/}）与真实工业订单流（WEPAStacks）两条互补主线开展系统研究与算法演进。在核心路径层，采用基于安全时间区间压缩的 SIPP 算法并结合增量预约表索引，显著降低了长时域避碰搜索开销；在决策层，面向一般分层与混合货架场景发展了拥堵关联大邻域修复（LNS）及高积压影子试派机制（BCRD），面向大规模真实订单流收敛出“事件门控滚动派单＋软分区联合遗传算法（GA）＋预约表 SIPP 物理预演”的综合方案（\code{gate}）。在涵盖 561 次独立验解运行的跨窗口多随机种子实验中，所有输出轨迹均 100\% 通过独立物理校验；在真实订单流既有验证窗口上，最终方案 \code{gate} 相对先到先服务（\code{fcfs}）、最早到达优先（\code{eta}）与滚动匈牙利指派（\code{hungarian}）分别将综合评价指标降低 5.45\%、2.43\% 与 3.87\%，并通过事件门控将开发窗 GA 触发次数削减 63\%。本文系统阐述了问题背景、方法框架与思路演进脉络，并为后续形式化建模、算法细节剖析与完整实验论证奠定结构基础。

    \vspace{5pt}
    \noindent{\sffamily\bfseries 关键词：}
    穿梭无人车（AGV）；多智能体路径规划（MAPF）；安全区间路径规划（SIPP）；联合遗传算法；事件门控滚动调度；立体仓库仿真
  \end{minipage}
\end{center}
\vspace{6pt}
\noindent\rule{\textwidth}{0.4pt}
\vspace{4pt}


% ==================== 已完成初稿部分（第 1–3 章） ====================
% ============================================================================
% 第 1 章 引言与问题背景
% ============================================================================
\section{引言与问题背景}
\label{sec:introduction}

\subsection{工业背景与研究意义}

随着智能制造、电商履约与柔性供应链的快速发展，传统人工拣选仓储模式在吞吐效率、空间利用率与人力成本等方面逐渐逼近瓶颈，自动化立体仓库（Automated Storage and Retrieval System, AS/RS）及多穿梭无人车（Automated Guided Vehicle, AGV）集群已成为现代智慧物流的核心基础设施。在典型的零件与货品立体仓库中，各类物料按批次和库位分布于密集排列的多层货架内；当上游生产系统或订单管理系统持续下发取货、出库与补货请求时，调度中心需要实时指挥多辆自主运行的 AGV 在有限通道空间内完成“空驶就位—货架取货—送达交接站—待命或最终返库”的闭环作业。

相比单车路径规划或离线静态排程，多取货任务并发下的仓储 AGV 调度在理论与工程上均属于典型的 NP-hard 组合优化与多智能体协同决策问题~\cite{ma2017lifelong, stern2019mapf}。在密集货架与狭窄巷道布局中，简单套用“最近车辆就近派单”或“静态最短路独立寻路”往往会引发严重的连锁退化：多辆车在同一交接站出口或主通道汇流形成拥堵与死锁，局部最优的即时派单破坏后续波次的任务均衡，进而拉长整体完工时间（Makespan）与订单响应延迟。因此，如何在严格满足物理避碰与仓储工艺规则的前提下，实现高吞吐、低延迟、低能耗且具备良好计算实时性的多 AGV 智能调度，具有直接的工程价值与重要的算法研究意义。

\subsection{多 AGV 并发调度的核心技术难点}

结合现代立体仓库的实际运行机理，多任务并发下的 AGV 智能调度主要面临以下四个相互交织的核心难点：

\begin{enumerate}
  \item \textbf{任务指派与时空路径规划的强耦合性}。车辆执行某项取货任务的真实代价并非静态几何距离，而是取决于当前在途车队的时空占位分布与巷道拥堵状态。若将任务分配与路径规划完全割裂，上层基于曼哈顿距离或静态最短路估算的“最优指派”在底层路径展开时可能遭遇大量等待与绕行；反之，若将全部车辆的离散时空状态联合纳入单一全局求解器，搜索空间随车辆数、订单数与规划时域呈指数级膨胀，难以满足在线滚动决策的时间约束~\cite{li2021rhcr}。
  \item \textbf{多层立体货架与复杂作业规则的硬约束}。真实零件仓库并非理想二维空旷栅格，而是包含单面多层货架、背靠背混合货架、单向/共享窄巷道、立柱障碍以及容量受限的交付工作站。任务不仅具有位置属性，还携带货架层高（直接决定升降装货耗时）及同面货架的层序解锁关系——同一货架面中必须遵循既定取货顺序，且交付站卸货后须及时驶离以释放终端资源。此外，入库与出库流可能部分重叠，空闲车辆在波次间隙需合理待命，并在订单流显式关闭后安全返库。
  \item \textbf{在线动态到达与因果信息边界}。在持续运行的真实订单流中，调度器无法预知未来尚未发布的订单，必须在严格遵守因果信息边界的前提下开展在线滚动决策。高负载高峰期的瞬时积压与低负载平峰期的零散到达交替出现，要求调度策略在“即时响应”与“适度协同筹划”之间取得动态平衡，同时避免因过度拟合特定历史波次而在未知订单窗口上产生泛化退化。
  \item \textbf{多目标权衡与物理可行性保障}。仓储运维通常同时关注三类相互制约的性能指标：决定系统整体周期效率的\emph{最终完工时间}（$T$）、反映单笔订单服务质量的\emph{平均响应耗时}（$R$），以及关联设备磨损与能耗的\emph{总行驶里程}（$D$）。此外，任何优化算法产生的调度方案都必须保证 100\% 满足顶点无碰撞、边交换无冲突、连续装卸时长与终端资源独占等底层物理约束，否则在物理系统中无法落地执行。
\end{enumerate}

\subsection{研究目标与本文主要贡献}

针对上述挑战，本项目以“理论建模清晰、算法演进有据、工程实现可复现”为原则，构建了面向多取货任务并发的仓库 AGV 智能调度与仿真验证体系。围绕课题要求，本文的主要研究工作与贡献可概括如下：

\begin{itemize}
  \item \textbf{构建分层解耦且严格可验的调度仿真架构}：设计并实现了纯 Python 标准库（零第三方重型依赖）的离散时空仿真执行内核与独立物理验解器，打通了“多工况建模—在线调度—冲突校验—可视化回放”的全链路，并支持交互式 Web 动态仿真回放与一键复现。
  \item \textbf{形成“通用多工况引擎”与“真实工业订单流”双主线协同研究}：一方面在自研通用引擎（\code{src/agv/}）中系统支持 Legacy 规则栅格、V4 分层单面货架及 V9/V9.1 混合背靠背货架布局，提出基于安全时间区间路径规划（SIPP）、拥堵关联大邻域修复（LNS）与高积压影子试派（BCRD）的模块化算法族；另一方面接入 SLAPStack/WEPAStacks 公开真实入出库订单流~\cite{pfrommer2022slapstack}，在 $62\times 47$ 大规模立体仓库网格（6 车、5 层、900\,s 订单窗口）下开展了多轮对照与跨窗口稳健性研究。
  \item \textbf{收敛出高效稳健的最终调度方案 \code{gate}}：在统一固定尺度线性目标函数下，通过问题驱动的多轮消融与跨窗口复评，最终定型“事件门控滚动派单＋软分区联合遗传算法（GA）＋预约表 SIPP 物理预演”方案（\code{gate}）。在涵盖 561 次独立验解运行的系统实验中，\code{gate} 在全部既有验证窗口上稳定优于 \code{fcfs}、\code{eta} 与 \code{hungarian} 基线，并通过事件门控大幅压缩低负载场景的无效搜索触发。
\end{itemize}

\subsection{本文结构安排}

本文其余部分组织如下：第~\ref{sec:methodology} 章给出系统总体架构、三层协同机制与最终采用方案概要；第~\ref{sec:evolution} 章沿问题驱动的视角系统梳理两条主线的思路推进、算法演化与关键消融决策；第~\ref{sec:modeling} 章阐述仓库几何、物理约束与评价指标的形式化数学模型；第~\ref{sec:algorithms} 章详述核心调度与路径规划算法设计；第~\ref{sec:experiments} 章展示跨场景、多随机种子及超参数敏感性实验结果与对比分析；第~\ref{sec:conclusion} 章总结全文工作并展望后续拓展方向。

% ============================================================================
% 第 2 章 解决方法概要
% ============================================================================
\section{解决方法概要}
\label{sec:methodology}

\subsection{三层协同调度架构}

为破解任务指派与时空路径规划之间的强耦合矛盾，同时兼顾大规模仓储环境下的计算效率与物理安全性，本系统采用\textbf{“上层任务指派—中层时空避碰—底层独立验解”}的三层协同架构（如图~\ref{fig:system-arch} 所示）。整个在线调度循环严格遵循因果信息边界：在每个离散时间步，调度器仅接收当前已发布且未完成的可见订单、车辆实时位姿及已承诺的未来时空预约表，输出原子可执行的调度计划并交由统一仿真内核步进推进。

\begin{figure}[htbp]
  \centering
  \begin{tikzpicture}[
    node distance=0.65cm and 0.9cm,
    box/.style={
      rectangle, rounded corners=3pt, draw=PrimaryNavy, thick,
      align=center, font=\small, inner sep=6pt, minimum height=1.15cm
    },
    layer/.style={
      rectangle, rounded corners=4pt, draw=BorderGray, dashed, thick,
      inner sep=8pt
    },
    arr/.style={-{Stealth[length=2.4mm]}, thick, color=PrimaryNavy},
    darr/.style={-{Stealth[length=2.2mm]}, thick, dashed, color=AccentTeal}
  ]
    % 顶部输入
    \node[box, fill=SoftBlueBg, minimum width=3.8cm] (env) {
      \sffamily\bfseries 仓库环境与动态订单流 \\
      \footnotesize 布局几何 / 层序状态 / 已发布订单
    };

    % 三层核心
    \node[box, fill=white, draw=AccentBlue, below=0.85cm of env, minimum width=4.3cm] (dispatch) {
      \sffamily\bfseries\color{PrimaryNavy} 1. 任务指派与决策层 \\
      \footnotesize 贪心 / 匹配 / 影子试派（BCRD）\\
      \footnotesize 事件门控 + 软分区联合 GA（\code{gate}）
    };

    \node[box, fill=white, draw=AccentTeal, below=0.85cm of dispatch, minimum width=4.3cm] (routing) {
      \sffamily\bfseries\color{AccentTeal} 2. 时空无碰撞路径规划层 \\
      \footnotesize 增量时空预约表 + SIPP 安全区间规划 \\
      \footnotesize 物理预演评分 / 拥堵关联 LNS 联合修复
    };

    \node[box, fill=SoftTealBg, draw=AccentTeal, below=0.85cm of routing, minimum width=4.3cm] (commit) {
      \sffamily\bfseries\color{PrimaryNavy} 3. 唯一仿真状态机与原子提交 \\
      \footnotesize 全未来冲突审计 $\rightarrow$ 步进推进 $\rightarrow$ 返库收尾
    };

    % 右侧独立验解与评估
    \node[box, fill=SoftAmberBg, draw=WarmAmber, right=1.35cm of routing, minimum width=3.5cm] (verify) {
      \sffamily\bfseries\color{WarmAmber} 独立物理验解器 \\
      \footnotesize 顶点/换边无碰撞校验 \\
      \footnotesize 装卸时长/层序/返库校验
    };

    \node[box, fill=SoftBlueBg, draw=PrimaryNavy, below=0.85cm of verify, minimum width=3.5cm] (metrics) {
      \sffamily\bfseries 综合指标与可视化 \\
      \footnotesize 指标 $T, R, D, F$ 统计 \\
      \footnotesize Web 回放 / 敏感性分析
    };

    % 连线
    \draw[arr] (env) -- node[right, font=\footnotesize, color=MutedText]{因果可见观测} (dispatch);
    \draw[arr] (dispatch) -- node[right, font=\footnotesize, color=MutedText]{候选任务分配与序列} (routing);
    \draw[arr] (routing) -- node[right, font=\footnotesize, color=MutedText]{无碰撞时空轨迹} (commit);

    % 反馈闭环（左侧）
    \draw[arr] (commit.west) -- ++(-1.15, 0) |- node[pos=0.25, left, font=\footnotesize, align=right, color=MutedText]{状态更新\\预约表滚动} (env.west);

    % 规划层与指派层双向预演反馈
    \draw[darr] (routing.east) -- ++(0.45, 0) |- node[pos=0.25, right, font=\footnotesize, color=AccentTeal]{预演代价} (dispatch.east);

    % 输出到独立验解
    \draw[arr] (commit.east) -| ($(verify.west)+(-0.35,0)$) |- (verify.west);
    \draw[arr] (verify) -- node[right, font=\footnotesize, color=MutedText]{100\% 可行确认} (metrics);
  \end{tikzpicture}
  \caption{多 AGV 仓库智能调度系统三层协同与独立验解架构图}
  \label{fig:system-arch}
\end{figure}

三层架构的核心职责划分如下：

\begin{itemize}
  \item \textbf{任务指派与决策层}：决定“由哪辆车以何种顺序响应哪些就绪任务”。在低负载、少量并发场景下采用低开销精确或贪心策略直接派单；在高负载或强竞争场景下，通过滚动影子试派（主线 A）或事件门控触发的软分区联合遗传搜索（主线 B）在多车多单组合空间内寻找全局协调解。
  \item \textbf{时空无碰撞路径规划层}：维护所有已指派车辆的离散时空预约表（记录顶点占用与有向边占用），并基于\textbf{安全区间路径规划（SIPP）}~\cite{phillips2011sipp} 将连续空闲时间段压缩为最大安全区间，从而在显著缩减状态搜索空间的同时，直接处理货架连续装卸占位、交付站即卸即走及窄巷道避让。此外，规划层一方面为上层候选方案提供真实的“预约表物理预演代价”，另一方面在主线 A 中支持对局部拥堵车队执行大邻域搜索（LNS）联合后缀重排~\cite{li2021mapflns}。
  \item \textbf{唯一仿真状态机与独立物理验解层}：所有规划结果在提交前均须通过全未来路径冲突与服务连续性审计，采用原子提交方式更新系统状态；仿真结束后，由与调度算法完全解耦的独立验解模块逐帧核查顶点冲突、相向换边冲突、装卸持续时长、同面货架层序解锁以及最终返库状态，确保实验数据绝对可信。
\end{itemize}

\subsection{双主线研究布局与最终采用方案}

考虑到课程项目既需要具备清晰通用、可扩展交互演示能力的工程平台，又需要在贴近工业界真实到达分布的大规模数据上验证核心算法的泛化性能，我们将整体研发组织为两条互补的研究主线，并明确了各自的最终定型版本与代码入口（见表~\ref{tab:dual-track}）：

\begin{table}[htbp]
  \centering
  \caption{本项目双研究主线定位、最终采用方案与核心机制总览}
  \label{tab:dual-track}
  \small
  \renewcommand{\arraystretch}{1.28}
  \begin{tabularx}{\textwidth}{>{\raggedright\arraybackslash}p{2.35cm} >{\raggedright\arraybackslash}p{3.0cm} >{\raggedright\arraybackslash}X >{\raggedright\arraybackslash}p{3.35cm}}
    \toprule
    \sffamily\bfseries 研究主线 & \sffamily\bfseries 最终采用方案 & \sffamily\bfseries 核心算法与关键机制 & \sffamily\bfseries 源码模块 \\
    \midrule
    \textbf{主线 B}：真实工业订单流（核心实验主线） &
    WEPA 推荐版 \code{gate}（算力受限时可选 \code{gate\_preview3}） &
    面向 $62\times 47$ 地图、6 车、5 层、900\,s 真实订单窗口：采用事件门控滚动派单（低负载精确分支/免触发 GA）、软分区联合 GA 寻优，并结合预约表 SIPP 物理预演择优提交。 &
    \footnotesize\code{research/}\newline\code{recommended-20261024/} \\
    \midrule
    \textbf{主线 A}：自研通用调度与可视化平台 &
    \code{src/agv/} v0.6.0（V9.1 + SIPP/BCRD 引擎） &
    支持 Legacy、V4 分层单面与 V9/V9.1 混合背靠背货架：V4 分层动态采用 \code{bcrd\_budget}（影子试派）；Legacy 静态采用 \code{sipp\_batch}；V9.1 静态采用 \code{sipp\_stream\_bcrd}（\code{max\_batches=2}）；内置本地 Web 回放。 &
    \footnotesize\code{src/agv/}（CLI 入口 \code{agv} 与 \code{agv serve}） \\
    \bottomrule
  \end{tabularx}
\end{table}

两条主线在底层共享同一套离散时空无碰撞建模思想与 SIPP 路径核心，但在实验侧重上形成良好互补：主线 A 侧重\textbf{多工况抽象与机理验证}，覆盖从规则栅格到单面多层、背靠背混合货架及流式关闭的渐进建模，并支持教师与评测人员通过浏览器直观观察 AGV 取货、避碰与返库过程；主线 B 则聚焦\textbf{真实工业数据驱动的深度算法优化}，依托公开的 WEPAStacks 真实入出库请求流（411,830 条原始记录抽样窗口）开展多方案配对比较、开发/验证集分离泛化测试与多随机种子敏感性分析，最终收敛出大作业核心推荐方案 \code{gate}。

\subsection{统一评价口径与工程复现保障}

为避免单一指标失衡，两条主线均采用兼顾周期效率、单单响应与能耗里程的加权综合目标函数 $F$（数值越低越优）：
\begin{itemize}
  \item \textbf{主线 B（独立物理尺度线性评分）}：去除对特定基线算法分母的依赖，按 900\,s 窗口时长与 6 车总步长尺度固定归一化，定义为
  \begin{equation}
    F_{\text{WEPA}} = 0.5\frac{T}{900} + 0.3\frac{R}{900} + 0.2\frac{D}{5400}
    \label{eq:obj-wepa-overview}
  \end{equation}
  其中 $T$ 为最后一笔订单的完成时刻（不含最终返库阶段），$R$ 为全部订单从发布到完成的平均耗时，$D$ 为含最终返库在内的全车队实际移动总步数。
  \item \textbf{主线 A（同场景参照归一化评分）}：在 V4 与 V9.1 等多尺度合成场景下，以同场景固定时空 A* 贪心基线指标 $(T_{\mathrm{ref}}, R_{\mathrm{ref}}, D_{\mathrm{ref}})$ 为参照基准 $1.0$，定义为 $F_{\text{Core}} = 0.5(T/T_{\mathrm{ref}}) + 0.3(R/R_{\mathrm{ref}}) + 0.2(D/D_{\mathrm{ref}})$；在 Legacy 场景下则直接对比平均流程时间 $R$、完工时间 $T$ 与总步数 $D$。
\end{itemize}

在工程实现方面，两条主线的核心调度、仿真与验解模块均仅依赖 \textbf{Python 3.10+ 标准库}，无需安装商业优化求解器或重型深度学习框架即可在跨平台环境下直接运行。同时，仓库内固化了完整的单元与集成测试套件、冻结证据校验脚本以及可重复的命令行复现入口，为后续章节的定量实验提供了严格的可审计基础。

% ============================================================================
% 第 3 章 思路推进与方案演化
% ============================================================================
\section{思路推进与方案演化}
\label{sec:evolution}

本项目的算法设计并非一次性套用现成模板，而是遵循\textbf{“构建基线 $\rightarrow$ 诊断瓶颈 $\rightarrow$ 提出改进 $\rightarrow$ 配对消融 $\rightarrow$ 跨窗口复评 $\rightarrow$ 架构收敛”}的实证迭代路径逐步推进。本章系统梳理两条研究主线的演化脉络，重点阐明每一阶段引入新机制的动因，以及若干探索性分支在严格跨窗口检验中被保留或舍弃的依据。

\subsection{总体演化路线总览}

图~\ref{fig:evolution-roadmap} 展示了两条主线的阶段演进关系。总体来看，项目前期通过主线 A 完成仓储物理建模的逐层细化（Legacy $\rightarrow$ V4 分层货架 $\rightarrow$ V9/V9.1 混合货架）与底层路径规划引擎的提速升级（时空 A* $\rightarrow$ SIPP $\rightarrow$ LNS/BCRD）；随后将经过验证的 SIPP 无碰撞执行内核引入主线 B，接入真实工业订单流 WEPAStacks，在统一线性目标口径下依次探索贪心规则、精确指派、群智能与遗传搜索、门控机制、自适应算子及深度强化学习调度切换，最终经由跨窗口多随机种子稳健性实验完成方案定型。

\begin{figure}[htbp]
  \centering
  \begin{tikzpicture}[
    node distance=0.45cm and 0.48cm,
    stageA/.style={
      rectangle, rounded corners=3pt, draw=AccentBlue, fill=SoftBlueBg,
      align=center, font=\footnotesize, inner sep=4.5pt, minimum width=2.45cm, minimum height=1.05cm
    },
    stageB/.style={
      rectangle, rounded corners=3pt, draw=AccentTeal, fill=SoftTealBg,
      align=center, font=\footnotesize, inner sep=4.5pt, minimum width=2.45cm, minimum height=1.05cm
    },
    finalB/.style={
      rectangle, rounded corners=3pt, draw=WarmAmber, fill=SoftAmberBg, very thick,
      align=center, font=\footnotesize\bfseries, inner sep=4.5pt, minimum width=2.55cm, minimum height=1.05cm
    },
    explore/.style={
      rectangle, rounded corners=3pt, draw=BorderGray, dashed, fill=white,
      align=center, font=\scriptsize, inner sep=4pt, minimum width=2.35cm, minimum height=0.85cm
    },
    arr/.style={-{Stealth[length=2mm]}, thick, color=PrimaryNavy},
    darr/.style={-{Stealth[length=1.8mm]}, dashed, color=MutedText}
  ]
    % 主线 A 标识
    \node[font=\small\sffamily\bfseries\color{AccentBlue}, anchor=west] (labA) at (0, 2.55) {主线 A：自研通用调度引擎演进（\code{src/agv/}）};

    \node[stageA] (a1) at (1.35, 1.55) {\bfseries Stage A1\\[-1pt] \code{baseline}\\ \scriptsize 贪心 + 时空 A*};
    \node[stageA, right=of a1] (a2) {\bfseries Stage A2\\[-1pt] \code{sipp\_budget/batch}\\ \scriptsize SIPP + LNS/批滚动};
    \node[stageA, right=of a2] (a3) {\bfseries Stage A3 (V4)\\[-1pt] \code{bcrd\_budget}\\ \scriptsize 分层货架 + 影子试派};
    \node[stageA, right=of a3] (a4) {\bfseries Stage A4 (V9/V9.1)\\[-1pt] \code{sipp\_stream\_bcrd}\\ \scriptsize 混合货架 + 流式关闭};

    \draw[arr] (a1) -- (a2);
    \draw[arr] (a2) -- (a3);
    \draw[arr] (a3) -- (a4);

    % 主线 B 标识（右移留出左侧连线通道，避免与箭线重叠）
    \node[font=\small\sffamily\bfseries\color{AccentTeal}, anchor=west] (labB) at (2.85, -0.35) {主线 B：WEPA 真实订单流算法收敛（\code{research/recommended-20261024/}）};

    \node[stageB] (b1) at (1.35, -1.35) {\bfseries Stage B1\\[-1pt] 真实流接入\\ \scriptsize 线性目标 + \code{fcfs/eta}};
    \node[stageB, right=of b1] (b2) {\bfseries Stage B2\\[-1pt] \code{simple} $\rightarrow$ \code{wide}\\ \scriptsize 软分区联合 GA};
    \node[finalB, right=of b2] (b3) {\bfseries Stage B3（最终）\\[-1pt] \code{gate}\\ \scriptsize 事件门控 + SIPP 预演};
    \node[stageB, right=of b3] (b4) {\bfseries Stage B4\\[-1pt] 跨窗口多 Seed 复评\\ \scriptsize 561 运行独立验解};

    \draw[arr] (b1) -- (b2);
    \draw[arr] (b2) -- (b3);
    \draw[arr] (b3) -- (b4);

    % SIPP 内核迁移连线（位于空白通道内，无文字重叠）
    \draw[darr, color=AccentTeal] (a2.south) -- ++(0, -0.42) -| node[pos=0.22, above, font=\scriptsize, color=AccentTeal]{SIPP 预约内核复用} (b1.north);

    % 探索分支对照
    \node[explore, below=0.55cm of b2] (e1) {\code{hungarian} / \code{hga}\\ \tiny 匈牙利匹配与混合初始化};
    \node[explore, below=0.55cm of b3] (e2) {\code{wide\_adapt} / 预算扫描\\ \tiny 自适应算子与 \code{preview3/gen24}};
    \node[explore, below=0.55cm of b4] (e3) {历史探索分支\\ \tiny 耦合修复 / \code{Double DQN}};

    \draw[darr] (b2.south) -- (e1.north);
    \draw[darr] (b3.south) -- (e2.north);
    \draw[darr] (b4.south) -- (e3.north);
  \end{tikzpicture}
  \caption{双主线问题驱动演化路线与探索性对照分支图}
  \label{fig:evolution-roadmap}
\end{figure}

\subsection{主线 A 演进：从基础启发式到分层动态与混合货架引擎}

主线 A 的演进以“先打通可验证基线，再逐级消除路径瓶颈、分配近视与场景失真”为核心逻辑，历经四个关键阶段：

\begin{enumerate}
  \item \textbf{阶段 A1（基线构建与瓶颈暴露）}：初始版本 \code{baseline} 采用距离贪心派单与标准时空扩展 A* 路径规划。随着任务密度和车辆数增加，时空 A* 逐时间步展开 $(x, y, t)$ 状态空间的开销迅速膨胀，且先到先占的顺序规划导致后规划车辆频繁在窄巷道内被动绕行或长时等待。
  \item \textbf{阶段 A2（SIPP 安全区间加速与局部协同）}：为压缩搜索空间，规划层升级为安全区间路径规划（SIPP）~\cite{phillips2011sipp}，将同一栅格连续无冲突的时间窗口合并为安全区间状态，并配套增量索引式时空预约表以降低冲突查询开销。在此基础上，针对局部拥堵引入大邻域搜索联合后缀重排（\code{joint\_only} / \code{sipp\_budget}），通过对高冲突车辆子集重新定序规划消除阻塞；针对静态并发场景进一步提出首轮多车有序批次滚动策略（\code{sipp\_batch}），在 Legacy 静态基准上将平均流程时间 $R$ 相对基线降低 2.62\%、总行驶步数减少 4.24\%，且运行速度明显优于同等质量的时空 A* 联合批规划（10.27\,s 对比 14.15\,s）。
  \item \textbf{阶段 A3（V4 分层单面货架与 BCRD 影子试派）}：为更贴合立体仓库物理工艺，V4 建模引入任务货架层高 \code{shelf\_level}、与层高成正比的装货耗时（$2\times\text{shelf\_level}$ 步）、交付站卸货即离以及同面货架层序解锁约束。在动态高积压工况下，单步即时指派容易因忽视后续待派任务而诱发货架面排队阻塞。为此提出\textbf{影子试派批次优化（BCRD，\code{bcrd\_budget}）}：当可见积压超过阈值时，在独立轻量影子环境中对多个候选有序批次序列开展前瞻滚动试派，择优提交当前首步指派。在 V4 分层动态 5-seed 实验中，\code{bcrd\_budget} 相对 \code{baseline} 将综合指标 $F$ 改善 3.01\%（4 胜 1 平 0 负），且平均单次求解耗时仅 0.17\,s。
  \item \textbf{阶段 A4（V9/V9.1 混合背靠背货架与流式关闭机制）}：进一步拓展至包含单排单层（双侧可取货）与背靠背多层货架的混合布局，显式支持取货侧选择、流式订单显式关闭事件（\code{close}）以及波次间隙的安全通道待命定位。在 V9.1 静态并发场景下，流式匹配与影子试派结合方案（\code{sipp\_stream\_bcrd}）将综合指标 $F$ 改善 1.53\%，且超参扫描表明批次上限 \code{max\_batches} 取 2 即可获得完整收益。
\end{enumerate}

\subsection{主线 B 演进：面向 WEPA 真实订单流的算法收敛}

为检验调度算法在大规模、非平稳真实工业订单流中的实战能力，主线 B 引入公开的 SLAPStack/WEPAStacks 真实入出库请求数据~\cite{pfrommer2022slapstack}，映射至 $62\times 47$ 大型仓库网格（6 辆 AGV、5 层货架、900\,s 滑动订单窗口），并围绕核心决策算法开展了四轮递进式演进与严格筛选：

\begin{enumerate}
  \item \textbf{评价重构与规则基线建立}：早期探索阶段（10 月 22 日）在合成数据上尝试了多类启发式与蚁群（ACO）、粒子群（PSO）路径混合，发现路径层引入元启发式对离散网格避碰收益有限，而真正的优化空间位于“订单分配与执行次序的联合决策”。接入真实入出库流后，首先重构了去基线依赖的固定尺度线性目标函数（式~\eqref{eq:obj-wepa-overview}），避免因参照基线波动导致跨轮评价失真；同时建立先到先服务（\code{fcfs}）、最早到达优先（\code{eta}）与滚动匈牙利指派（\code{hungarian}）~\cite{kuhn1955hungarian} 三类规则与局部匹配基线。
  \item \textbf{从单轮贪心到软分区联合 GA（\code{simple} $\rightarrow$ \code{wide}）}：规则基线与一次性匈牙利匹配仅优化当前瞬时批次，难以协调多辆 AGV 连续多单的任务均衡与交付站错峰。为此引入软分区联合遗传算法，将未完成订单的车辆归属与执行次序联合编码，并由轻量模型评估与预约表 SIPP 物理预演联合打分。从初期受限预算版 \code{simple} 扩展种群与迭代预算至 \code{wide} 后，调度质量获得显著提升，但也暴露出在低负载窗口高频触发 GA 带来的计算冗余问题。
  \item \textbf{事件门控与最终方案定型（\code{gate}）}：针对 \code{wide} 的算力冗余，在 \code{gate} 中引入\textbf{事件门控滚动派单}机制：当在途与就绪任务数处于低负载区间时，跳过遗传迭代，直接调用位掩码动态规划或低开销精确分支完成分配；仅当多车多单存在显著竞争时才触发软分区联合 GA，并对顶级候选执行真实预约表 SIPP 物理预演择优。多随机种子对照表明，\code{gate} 与 \code{wide} 在综合目标 $F$ 上高度接近（开发窗 +0.11\%，验证窗 $-0.10\%$），同时将开发窗的 GA 触发次数由 560 次降至 205 次（削减 63\%），且在 4 个既有验证窗口上相对 \code{fcfs}、\code{eta}、\code{hungarian} 分别取得 +5.45\%、+2.43\% 与 +3.87\% 的全胜改善（4/0/0）。此外，预算扫描进一步确认：将完整 SIPP 预演数降至 3 次（\code{gate\_preview3}）可在 $F$ 仅波动 $-0.33\%$ 的代价下削减约 44\% 的 CPU 耗时，形成面向算力受限场景的实用轻量档。
  \item \textbf{探索性分支的跨窗口稳健性复盘}：在选型收敛过程中，项目秉持“开发集探索、独立验证集冻结复评”的科研规范，系统检验了多项进阶改进设想：
  \begin{itemize}
    \item \emph{匈牙利混合初始化（\code{hga}）与自适应交叉/变异（\code{wide\_adaptXY}）}：在开发窗上，\code{hga} 相对 \code{gate} 取得 +0.95\% 改善，仅开启交叉自适应的 \code{wide\_adapt01} 相对固定参数版取得 +0.94\% 改善；然而在 4 个独立既有验证窗口复评中，两者均出现跨窗回退（\code{hga} 为 $-1.12\%$、0/0/4；\code{wide\_adapt01} 为 $-0.54\%$、1/0/3），说明局部开发窗上的单点增益存在过拟合风险。
    \item \emph{扩大迭代代数与事后局部修复}：将 GA 代数由 16 增至 24（\code{gate\_gen24}）在验证窗反而下降 1.41\%（0/0/4），表明在动态滚动环境中过度加深单步静态搜索会加剧对当前不完全信息的过拟合；开放流终点预测、瓶颈邻域修复（ALNS 思路）~\cite{ropke2006alns} 与双车耦合局部重排在开发/验证集间亦存在权衡不一致，故保留为对照算子而非默认主线。
    \item \emph{深度强化学习切换探索（\code{Double DQN}）}：在独立归档分支 \code{dqn-experiment-v2} 中探索了以 Double DQN~\cite{vanhasselt2016ddqn} 动态选择局部匹配与自适应 GA 动作的元调度框架；测试结果显示其整体综合目标略低于纯自适应 GA 基线（$-0.37\%$），且引入额外训练依赖，因此予以完整归档供对比反思，未纳入零依赖默认主线。
  \end{itemize}
\end{enumerate}

\subsection{关键演进阶段与算法对照汇总}

为直观呈现上述思路推进脉络，表~\ref{tab:evolution-summary} 汇总了两条主线从初始基线到最终方案的关键演进节点、所要解决的核心痛点及验证定位。

\begin{table}[htbp]
  \centering
  \caption{核心算法演进历程、解决痛点与最终定位对照表}
  \label{tab:evolution-summary}
  \footnotesize
  \renewcommand{\arraystretch}{1.24}
  \begin{tabularx}{\textwidth}{>{\raggedright\arraybackslash}p{2.75cm} c >{\raggedright\arraybackslash}p{2.85cm} >{\raggedright\arraybackslash}X >{\raggedright\arraybackslash}p{2.45cm}}
    \toprule
    \sffamily\bfseries 方案标识 & \sffamily\bfseries 主线 & \sffamily\bfseries 针对的核心痛点 & \sffamily\bfseries 关键技术突破 & \sffamily\bfseries 验证定位 \\
    \midrule
    \code{baseline} & A & 缺乏可复现对照参照 & 距离贪心派单 + 时空扩展 A* 路径规划 & 通用基线 \\
    \code{sipp\_budget} / \code{sipp\_batch} & A & 时空 A* 慢、局部巷道拥堵与静态首轮失衡 & SIPP 安全区间压缩 + 增量预约表 + 拥堵关联 LNS 与首轮批滚动 & Legacy 静态优选（$R$ 降 2.62\%） \\
    \code{bcrd\_budget} & A & V4 分层货架高积压下即时派单近视 & 轻量影子环境前瞻试派多候选有序批次序列 & V4 动态优选（$F$ 降 3.01\%） \\
    \code{sipp\_stream\_bcrd} & A & V9.1 混合背靠背货架与显式流关闭协同 & 显式取货侧选择 + 安全待命 + BCRD（\code{max\_batches=2}） & V9.1 静态优选（$F$ 降 1.53\%） \\
    \midrule
    \code{fcfs} / \code{eta} / \code{hungarian} & B & 真实订单流下缺少规则与一次性匹配参照 & 按到达、预计耗时或当前代价矩阵即时指派 & 规则与匹配基线 \\
    \code{simple} $\rightarrow$ \code{wide} & B & 单批次匹配忽视多单连续协同与错峰 & 软分区联合 GA 编码多车任务序列并扩大搜索预算 & GA 演进过渡版 \\
    \textbf{\code{gate}} & \textbf{B} & \code{wide} 在低负载频繁触发 GA 致算力冗余 & \textbf{事件门控低负载精确分支 + 软分区联合 GA + SIPP 物理预演} & \textbf{最终默认推荐}（验证窗全胜基线） \\
    \code{gate\_preview3} & B & 离线/边缘算力受限下的耗时压缩需求 & 完整 SIPP 预演数由默认降至 3 次，CPU 降约 44\% & 算力受限轻量档 \\
    \code{hga} / \code{wide\_adapt} / \code{DQN} & B & 探索初始化加速、自适应算子与 RL 元调度 & 匈牙利初始个体、自适应交叉/变异率、Double DQN 动作切换 & 探索对照与归档（跨窗未优于 \code{gate}） \\
    \bottomrule
  \end{tabularx}
\end{table}

通过上述两主线交织、正负对照并重的演化过程，本项目不仅确立了综合表现最稳健的默认方案 \code{gate} 及通用引擎 \code{src/agv/}，也为后续章节展开严格的数学建模、算法伪代码设计与全量多维实验分析构建了清晰的逻辑主干。


% ==================== 待与队友商议合写部分（第 4–7 章提纲） ====================
% ============================================================================
% 第 4 章 问题形式化建模与物理约束（待与队友商议合写）
% ============================================================================
\section{问题形式化建模与物理约束}
\label{sec:modeling}

\begin{planbox}{第 4 章写作规划（供双人分工确认）}
  本章拟以严谨的运筹学与多智能体路径规划（MAPF）记号给出全系统的统一数学表述，建议细分为以下小节并在商定分工后补充完整公式与示意配图：
  \begin{enumerate}
    \item \textbf{4.1 仓库拓扑图与多层货架建模}：定义四邻接离散栅格图 $G=(V,E)$、通道/货架/立柱/交付站/停靠区集合，分别给出主线 A（Legacy / V4 单面多层 / V9.1 混合背靠背货架）与主线 B（$62\times 47$、6 车、5 层 WEPAStacks 映射网格）的布局示意图。
    \item \textbf{4.2 订单任务与生命周期状态机}：形式化描述订单发布时刻 $r_j$、SKU 与货架层高 $l_j$、入/出库任务重叠规则、装货持续时间（如 V4 的 $2 l_j$ 步）、交付站卸货即离、波次间安全待命及流关闭（\code{close}）后的最终返库状态转移。
    \item \textbf{4.3 时空无碰撞与仓储工艺硬约束}：列出离散步长下顶点互斥约束、相向边交换互斥约束、连续装卸占位不变性、同面货架层序解锁偏序关系以及在线调度的因果信息边界条件。
    \item \textbf{4.4 评价指标与目标函数定义}：给出完工时间 $T$、平均响应时间 $R$、含返库总里程 $D$、主线 B 固定尺度线性目标 $F_{\text{WEPA}}$、主线 A 参照归一化目标 $F_{\text{Core}}$ 以及统一改善率公式的严格定义；并简要说明微型两车 CP-SAT 精确参照模型的作用。
  \end{enumerate}
\end{planbox}

% ============================================================================
% 第 5 章 核心调度与路径规划算法设计（待与队友商议合写）
% ============================================================================
\section{核心调度与路径规划算法设计}
\label{sec:algorithms}

\begin{planbox}{第 5 章写作规划（供双人分工确认）}
  本章重点剖析系统最终采用的核心算法机理（建议配以规范伪代码、流程图与时间复杂度分析，避免整段粘贴源码）：
  \begin{enumerate}
    \item \textbf{5.1 基于增量时空预约表的 SIPP 路径规划}：阐述如何将栅格上的动态障碍占位压缩为最大安全时间区间，支持单车“就位—取货连续占位—送货—离站—返库”组合路径拼接，以及增量预约表对搜索效率的加速作用~\cite{phillips2011sipp}。
    \item \textbf{5.2 面向真实订单流的事件门控软分区联合 GA（\code{gate}）}：
    \begin{itemize}
      \item 阐述事件门控触发条件与低负载位掩码动态规划分支；
      \item 详述软分区联合染色体编码、交叉/变异算子与多候选种群管理~\cite{vidal2022hgs}；
      \item 说明“快速代理评估 + Top-$K$ 候选真实预约表 SIPP 物理预演”的两级评分机制，以及 \code{gate\_preview3} 轻量档的预算权衡原理。
    \end{itemize}
    \item \textbf{5.3 通用引擎核心机制（LNS 联合修复与 BCRD 影子试派）}：介绍主线 A 中的拥堵关联车队邻域选择与 LNS 后缀重排（\code{sipp\_budget} / \code{sipp\_batch}）~\cite{li2021mapflns}，以及高积压工况下独立影子环境前瞻试派（\code{bcrd\_budget} / \code{sipp\_stream\_bcrd}）的算法流程。
    \item \textbf{5.4 工程架构与独立物理验解设计}：概述零第三方依赖标准库实现、原子提交审计、独立物理验解模块（\code{validation.py} / \code{verify.py}）以及本地 Web 交互式仿真可视化界面（\code{agv serve}）。
  \end{enumerate}
\end{planbox}

% ============================================================================
% 第 6 章 实验设计与结果分析（待与队友商议合写）
% ============================================================================
\section{实验设计与结果分析}
\label{sec:experiments}

\begin{planbox}{第 6 章写作规划（供双人分工确认）}
  本章将全面呈现 561 次独立验解运行的定量结果与可视化图表（原始数据与高清图片均已冻结归档于代码仓库 \code{results/} 目录，后续可直接转入本仓库 \code{figures/} 引用）：
  \begin{enumerate}
    \item \textbf{6.1 实验设置与复现协议}：说明硬件环境、软件零依赖复现命令、开发窗与既有验证窗的窗口划分、随机种子配置（WEPA seed 0--2 / Core 5 seeds）及 100\% 独立物理验解通过情况。
    \item \textbf{6.2 主线 B 核心对照实验（WEPA 开发窗与既有验证窗）}：
    \begin{itemize}
      \item 展示 \code{fcfs}、\code{eta}、\code{hungarian}、\code{simple}、\code{wide}、\code{gate} 与 \code{hga} 在 4 个开发窗（200 次运行）及 4 个既有验证窗（156 次运行）上的 $T, R, D, F$ 与配对胜平负统计表；
      \item 结合多 seed 分布箱线图与窗级配对对比图，论证 \code{gate} 相对规则/匹配基线的稳定优势（+5.45\% / +2.43\% / +3.87\%）以及相对 \code{wide} 削减 63\% GA 触发次数的效率收益；
      \item 分析 \code{hga} 与 \code{wide\_adapt01} 在开发窗领先但验证窗回退的跨窗泛化现象，以及微型两车 CP-SAT 精确参照的误差检验。
    \end{itemize}
    \item \textbf{6.3 超参数敏感性与算力—质量权衡分析}：结合 2$\times$2 预算扫描（\code{gate\_preview3} / \code{gate\_gen8} / \code{gate\_gen24}）与 2$\times$2 自适应算子扫描（\code{wide\_adapt00..11}），展示 CPU 耗时与 $F$ 值的帕累托权衡曲线，并简述历史 \code{Double DQN} 分支对照结果。
    \item \textbf{6.4 主线 A 多工况演进与参数扫描（Legacy / V4 / V9.1）}：汇总 205 次 5-seed 实验结果，对比 \code{bcrd\_budget}（V4 动态 +3.01\%）、\code{sipp\_batch}（Legacy 静态 $R$ 改善 +2.62\%）、\code{sipp\_stream\_bcrd}（V9.1 静态 +1.53\%）以及 \code{max\_batches} 与 \code{trial\_budget} 敏感性曲线，并附本地 Web 可视化仿真关键帧截图。
  \end{enumerate}
\end{planbox}

% ============================================================================
% 第 7 章 总结与展望
% ============================================================================
\section{总结与展望}
\label{sec:conclusion}

\subsection{全文工作总结}

本文面向自动化立体仓库中多取货任务并发下的多穿梭无人车（AGV）智能调度问题，围绕“任务指派”与“离散时空无碰撞路径规划”的耦合挑战，完成了从多层/混合货架形式化建模、核心算法演进、零依赖工程架构落地到大规模多随机种子实证检验的完整研究闭环。主要成果归纳如下：

\begin{enumerate}
  \item \textbf{构建了严格可验的三层协同架构与零依赖通用仿真平台（\code{src/agv/} v0.6.0）}：渐进建立了涵盖 Legacy 规则栅格、V4 单面多层货架（层高相关装货耗时与同面层序解锁）及 V9/V9.1 混合背靠背货架（显式选侧、开放流显式关闭与安全待命）的数学模型；底层通过增量索引式时空预约表、安全区间路径规划（SIPP）与线程局部几何缓存大幅压缩长时域避碰搜索开销；顶层实现解耦的独立物理验解器与本地交互式 Web 可视化界面（\code{agv serve}），确保所有实验输出 100\% 物理合法且支持一键复现。
  \item \textbf{针对多工况特征发展了 LNS 联合修复、BCRD 反事实影子试派与 P90 尾部保护机制}：在通用引擎主线 A 中，静态首轮批滚动（\code{sipp\_batch}）将 Legacy 静态流程时间改善 $+2.63\%$、总步数减少 $4.24\%$；高积压反事实影子试派（\code{bcrd\_budget}）在 V4 动态 5-seed 实验中取得 $+3.01\%$ 综合改善（单次耗时 $0.17\,\text{s}$），配合无损事件驱动影子推进在 90 个留出场景上削减 $43.2\%$ CPU 开销；在 V9.1 的 960 次未见种子留出实验中，历史热度安全选点与 ETA 匹配（\code{G4}）及流式 BCRD（\code{B0}/\code{B1}）分别在动态突发与密集积压工况下取得显著正收益，其中 P90 尾部护栏将留出集尾部退化例数从 13/240 清零至 0/240。
  \item \textbf{在 WEPAStacks 真实工业订单流上通过跨窗口多 Seed 复评确立了最终方案 \code{gate}}：在 $62\times 47$ 网格、6 车、5 层、900\,s 真实订单窗口及固定尺度线性目标下，经由 200 次开发窗与 156 次既有验证窗多随机种子配对实验证实，“事件门控滚动派单＋软分区联合 GA＋预约表 SIPP 物理预演”（\code{gate}）在全部 4 个既有验证窗口上稳定战胜 \code{fcfs}（$+5.45\%$）、\code{eta}（$+2.43\%$）与 \code{hungarian}（$+3.87\%$），并通过事件门控将开发窗 GA 触发次数削减 $63\%$；轻量档 \code{gate\_preview3} 可在 $F$ 仅变化 $-0.33\%$ 的条件下节省约 $44\%$ CPU 时间。同时，跨窗复评定量揭示了匈牙利初始注入（\code{hga}）、单开交叉自适应（\code{wide\_adapt01}）、变异自适应饱和（$>86\%$ 代数触顶 $0.95$）及过度扩大代数（\code{gate\_gen24}）在未知订单窗口上的过拟合回退机理，并完成了 Double DQN 元调度切换的对比归档。
\end{enumerate}

\begin{teammatebox}{7.1 节补充留白（可选 —— 供队友补充总结）}
  队友在完成 4.3、4.4、5.3、6.2 及 6.4.4 节内容后，可在此处补充关于 WEPA 真实流建模、\code{gate} 软分区编码设计及早期阶段消融实验的总结凝练（定稿时可直接合并入上方总结列表并移除本提示框）。
\end{teammatebox}

\subsection{局限性分析与未来展望}

尽管本系统在多类合成与真实工业订单窗口上展现出稳健的调度质量与工程可复现性，后续研究仍可沿以下方向进一步深化：
\begin{itemize}
  \item \textbf{更大规模车队与高拥堵窗口联合重排（RHCR 融合）}：当前实验主要聚焦 4--6 辆 AGV 与 3--5 层立体货架；当车队规模扩展至 12--20 辆且巷道密度进一步提高时，可结合滚动时域冲突消解（RHCR）~\cite{li2021rhcr} 在有限时间窗内引入周期性多车路径重规划。
  \item \textbf{残余尾估计引导的浅层影子试派}：BCRD 消融实验表明，若无可靠的残余完工时间估计，单纯截断影子推演深度易触发安全回退；未来可结合离线学习的价值函数拟合残余积压代价，在超大订单池下进一步压缩前瞻步数。
  \item \textbf{电量约束与多工作站动态协同}：在实际全天候仓储运行中，可进一步纳入 AGV 电池能耗模型、充电桩排队调度以及多拣选工作站动态负载均衡，推动算法向工业级全要素数字孪生系统演进。
\end{itemize}


% ==================== 参考文献 ====================
\bibliographystyle{unsrt}
\bibliography{references}

\end{document}
